\documentclass[../../main.tex]{subfiles}%
\graphicspath{{\subfix{./image/}}{\subfix{../../image/}}}

\begin{document}

\section{参考文献与注记}\label{section:ch02-05}

我们的主要依据同样是 Federer\cite{federer1969}，特别是\cite{federer1969} 第 2.10 节。Steiner 对称化可在\cite{federer1969} 第 2.10.30、2.10.31 节中找到。等直径不等式（\cref{theorem:ch02-2-4}）的证明我们紧密遵循 Hardt\cite{hardt1979}，但采纳了 L--F Tam 给出的一个简化：他指出只需沿坐标方向作对称化。R. Hardt 向我们介绍了 Tam 的这一观察。% XREF p=115 "for the proof of the isodiametric inequality"

\cref{theorem:ch02-2-5}（即 $\mathcal{H}^n=\mathcal{L}^n$）的证明来自 Hardt\cite{hardt1979} 与 Simon\cite{simon1984} 第 2.3--2.6 节。关于\cref{section:ch02-03}中的密度定理，我们参考了\cite{simon1984} 第 3 节。% XREF p=115 "The proof that H^n = L^n is from Hardt [H] and Simon [S, Sections 2.3–2.6]" | p=115 "for the density theorems in Section 2.3"

Falconer\cite{falconer1985,falconer1990} 与 Morgan\cite{morgan2009} 对 Hausdorff 测度给出了很好的入门介绍。一本很好的进阶教材是 Mattila\cite{mattila1995}。

\end{document}
